\documentclass[11pt]{article}
\usepackage[margin=0.8in]{geometry}
\usepackage{amsmath,amssymb,amsthm,microtype}
\newtheorem{theorem}{Theorem}
\begin{document}
\begin{center}
{\Large A Boundary for Bootstrapping from Convergent Base Cylinders}\par\medskip
Collatz proof project\quad---\quad September 5, 2026
\end{center}
The inverse-cycle construction supplies conditional transfer certificates on
$u=r+2^dQ$ once the two base partners $3r+2$ and $27r+20$ have reached the
trivial cycle. It can repair either cycle-charge mismatch, but its binary
precision must accommodate the base target's path to the cycle.
We test whether such certificates can bootstrap all parameters from strictly
smaller canonical bases. The answer is negative for an explicit infinite family.

\begin{theorem}[Necessary size condition]
For every natural $n,t$, $n\leq2^tT^t(n)$, where $T$ is the shortcut
Collatz map. Consequently,
\[
 T^d(27r+20)\leq2\quad\Longrightarrow\quad27r+20\leq2^{d+1}.
\]
\end{theorem}
\begin{proof}
The exact affine iterate formula is
$2^tT^t(n)=3^{j_t(n)}n+c_t(n)$, where $j_t(n)$ and $c_t(n)$ are
nonnegative integers. Its right side is at least $n$.
Apply this to the target and its assumed endpoint bound.
Equivalently, no shortcut step can reduce a positive value faster than halving.
\end{proof}

\begin{theorem}[Top residues cannot be cycle bases at their own depth]
For every natural $d$,
\[
 2<T^d\bigl(27(2^d-1)+20\bigr).
\]
\end{theorem}
\begin{proof}
The contrary would imply $27\cdot2^d-7\leq2\cdot2^d$.
Thus $25\cdot2^d\leq7$, impossible since $2^d\geq1$.
This argument covers every depth, not a finite search interval.
\end{proof}

\begin{theorem}[No smaller cycle-base certificate for Mersenne parameters]
Let $u=2^k-1$. There are no natural $d,r$ satisfying all four conditions
\[
 r<2^d,\qquad r=u\bmod2^d,\qquad r<u,\qquad
 T^d(27r+20)\leq2.
\]
\end{theorem}
\begin{proof}
If $d>k$, then $u<2^d$ and its canonical residue is $u$ itself,
contradicting $r<u$. Hence $d\leq k$.
Since $2^d$ divides $2^k$, the residue of $2^k-1$ is $2^d-1$.
The preceding theorem excludes the required target endpoint.
Lean checks the residue identity and both cases with natural-number arithmetic.
\end{proof}

\paragraph{Implication for the current strategy.}
For instance, proper binary prefixes of parameter 31 have values
1, 3, 7, and 15. None can serve as a base whose target has reached the
cycle within that prefix depth. Choosing more precision eventually selects
31 itself as the base, reintroducing the convergence obligation at that
same parameter. The theorem applies to every $2^k-1$, not only this example.

No choice of later positive or negative cycle excursions repairs this base
selection problem: the obstruction precedes those excursions. Increasing the
number of certificates constructed by first verifying their base targets does
not establish the missing global induction coverage.

\paragraph{Scope and next obligation.}
This is not a nonconvergence theorem and does not exclude direct merges before
the cycle, other residue conventions, other induction orders, or a separate
uniform treatment of these parameters. It excludes only the stated bootstrap
from proper smaller canonical bases that have reached the target cycle by the
binary depth. In particular, it does not contradict the inverse excursion's
exact transfer identities or its conditional local construction.

\paragraph{Verification.}
\begingroup\raggedright
Five declarations in \texttt{Collatz.CycleCylinderBoundary} prove the size bound,
the cycle-entry restriction, the top-residue inequality, the Mersenne residue
identity, and the final all-depth exclusion. They are included in the selected
Lean axiom audit. No mathematical priority claim is made. Collatz remains unproved.
\par\endgroup
\end{document}
